% GENERATED FILE. Edit canonical lessons, then run npm run generate:books.
% canonical-source-hash: fnv1a64:b2d87e9904f9a1a8
% canonical-lessons: RU-C460-ruchey, RU-C460-boloto, RU-C460-rodnik, RU-C460-prud, RU-C460-zhurchat

\chapter{Streams and ponds}
\label{ch:ru-streams-and-ponds}
\hlchaptermodality{460}

\begin{chapteropening}
\textbf{By the end of this chapter:} \emph{I can describe water in the countryside: a stream, a swamp, a spring, a pond and water babbling.}

Complete the last lesson of chapter 460: I can describe water in the countryside: a stream, a swamp, a spring, a pond and water babbling.
\end{chapteropening}

\section[\texorpdfstring{\ru{ручей} (ruchéy)}{ручей (ruchéy)}]{\ru{ручей} (ruchéy) — a stream, a brook}

\label{lesson:RU-C460-ruchey}



\begin{quote}
Before the new one: say the Russian for to scramble, then the Russian for to stumble.
\end{quote}

\subsection*{You'll want to know: \ru{ручей}}
\textbf{\ru{ручей}} (\emph{ruchéy}) — \textquotedblleft{}a stream, a brook\textquotedblright{}.

\begin{quote}
\ru{В} \ru{лесу} \ru{мы} \ru{пили} \ru{воду} \ru{из} \ru{ручья}.
\end{quote}

\begin{quote}
\emph{V lesú my píli vódu iz ruchyá.}
\end{quote}

\begin{quote}
In the forest we drank water from a stream.
\end{quote}

Masculine; the \ru{е} drops out: \emph{\ru{ручья}}, \emph{\ru{ручьи}}.

\begin{cousinweb}
\emph{\ru{ручей}} is inherited from Proto-Slavic \emph{ručajĭ} 'a loudly flowing stream' (EtymDB), which EtymDB builds on a verb \emph{ručati} 'to roar, to hum'.
\end{cousinweb}

\begin{grammarlens}[title={What you've built}]
One more word for telling, explaining and describing at B1.
\end{grammarlens}

\subsection*{Your turn}
\begin{itemize}
\raggedright
  \item \emph{Say it:} \emph{ruchéy}
  \item \emph{Say it:} \emph{ruchéy}, once more
  \item \emph{Say it:} a sentence of your own with \emph{ruchéy}
  \item \emph{Recall:} say the Russian for to scramble, then the Russian for to stumble, then say \emph{ruchéy} again
\end{itemize}

\begin{tcolorbox}[breakable,colback=teal!4,colframe=teal!35!black,title={Before you move on}]
What does \ru{ручей} mean? (\textquotedblleft{}A stream, a brook\textquotedblright{}.) Say it once more.
\end{tcolorbox}
\section[\texorpdfstring{\ru{болото} (bolóto)}{болото (bolóto)}]{\ru{болото} (bolóto) — a swamp, a marsh}

\label{lesson:RU-C460-boloto}



\begin{quote}
Before the new one: say the Russian for to stumble, then the Russian for a stream.
\end{quote}

\subsection*{You'll want to know: \ru{болото}}
\textbf{\ru{болото}} (\emph{bolóto}) — \textquotedblleft{}a swamp, a marsh\textquotedblright{}.

\begin{quote}
\ru{За} \ru{деревней} \ru{большое} \ru{болото}.
\end{quote}

\begin{quote}
\emph{Za derévney bolshóye bolóto.}
\end{quote}

\begin{quote}
Beyond the village there is a big swamp.
\end{quote}

Neuter. The native \emph{-\ru{оло}-} shape; Church Slavonic has \emph{\ru{блато}}.

\begin{cousinweb}
\emph{\ru{болото}} is inherited through Old East Slavic \emph{boloto} from Proto-Slavic \emph{bolto} 'a swamp' (EtymDB); its \emph{-\ru{оло}-} is the native East Slavic shape.
\end{cousinweb}

\begin{grammarlens}[title={What you've built}]
One more word for telling, explaining and describing at B1.
\end{grammarlens}

\subsection*{Your turn}
\begin{itemize}
\raggedright
  \item \emph{Say it:} \emph{bolóto}
  \item \emph{Say it:} \emph{bolóto}, once more
  \item \emph{Say it:} a sentence of your own with \emph{bolóto}
  \item \emph{Recall:} say the Russian for to stumble, then the Russian for a stream, then say \emph{bolóto} again
\end{itemize}

\begin{tcolorbox}[breakable,colback=teal!4,colframe=teal!35!black,title={Before you move on}]
What does \ru{болото} mean? (\textquotedblleft{}A swamp, a marsh\textquotedblright{}.) Say it once more.
\end{tcolorbox}
\section[\texorpdfstring{\ru{родник} (rodník)}{родник (rodník)}]{\ru{родник} (rodník) — a spring (of water)}

\label{lesson:RU-C460-rodnik}



\begin{quote}
Before the new one: say the Russian for a stream, then the Russian for a swamp.
\end{quote}

\subsection*{You'll want to know: \ru{родник}}
\textbf{\ru{родник}} (\emph{rodník}) — \textquotedblleft{}a spring (of water)\textquotedblright{}.

\begin{quote}
\ru{Вода} \ru{в} \ru{роднике} \ru{холодная} \ru{и} \ru{чистая}.
\end{quote}

\begin{quote}
\emph{Vodá v rodniké kholódnaya i chístaya.}
\end{quote}

\begin{quote}
The water in the spring is cold and clean.
\end{quote}

Masculine; the stress moves: \emph{\ru{родника}} (\emph{rodniká}).

\begin{cousinweb}
\emph{\ru{родник}} is usually derived from \emph{\ru{род}-} 'birth', as in \emph{\ru{родить}} 'to give birth': water born from the earth. EtymDB lists the word without a source, so the line is hedged.
\end{cousinweb}

\begin{grammarlens}[title={What you've built}]
One more word for telling, explaining and describing at B1.
\end{grammarlens}

\subsection*{Your turn}
\begin{itemize}
\raggedright
  \item \emph{Say it:} \emph{rodník}
  \item \emph{Say it:} \emph{rodník}, once more
  \item \emph{Say it:} a sentence of your own with \emph{rodník}
  \item \emph{Recall:} say the Russian for a stream, then the Russian for a swamp, then say \emph{rodník} again
\end{itemize}

\begin{tcolorbox}[breakable,colback=teal!4,colframe=teal!35!black,title={Before you move on}]
What does \ru{родник} mean? (\textquotedblleft{}A spring (of water)\textquotedblright{}.) Say it once more.
\end{tcolorbox}
\section[\texorpdfstring{\ru{пруд} (prud)}{пруд (prud)}]{\ru{пруд} (prud) — a pond}

\label{lesson:RU-C460-prud}



\begin{quote}
Before the new one: say the Russian for a swamp, then the Russian for a spring.
\end{quote}

\subsection*{You'll want to know: \ru{пруд}}
\textbf{\ru{пруд}} (\emph{prud}) — \textquotedblleft{}a pond\textquotedblright{}.

\begin{quote}
\ru{В} \ru{пруду} \ru{плавают} \ru{утки}.
\end{quote}

\begin{quote}
\emph{V prudú plávayut útki.}
\end{quote}

\begin{quote}
Ducks are swimming in the pond.
\end{quote}

Masculine; \emph{\ru{в} \ru{пруду}} — in the pond.

\begin{cousinweb}
\emph{\ru{пруд}} is usually counted an inherited Slavic word; EtymDB gives no source but lists Polish \emph{prąd} and Slovak \emph{prúd} 'a current' as its cousins.
\end{cousinweb}

\begin{grammarlens}[title={What you've built}]
One more word for telling, explaining and describing at B1.
\end{grammarlens}

\subsection*{Your turn}
\begin{itemize}
\raggedright
  \item \emph{Say it:} \emph{prud}
  \item \emph{Say it:} \emph{prud}, once more
  \item \emph{Say it:} a sentence of your own with \emph{prud}
  \item \emph{Recall:} say the Russian for a swamp, then the Russian for a spring, then say \emph{prud} again
\end{itemize}

\begin{tcolorbox}[breakable,colback=teal!4,colframe=teal!35!black,title={Before you move on}]
What does \ru{пруд} mean? (\textquotedblleft{}A pond\textquotedblright{}.) Say it once more.
\end{tcolorbox}
\section[\texorpdfstring{\ru{журчать} (zhurchát)}{журчать (zhurchát)}]{\ru{журчать} (zhurchát) — to babble, to murmur (of water)}

\label{lesson:RU-C460-zhurchat}



\begin{quote}
Before the new one: say the Russian for a spring, then the Russian for a pond.
\end{quote}

\subsection*{You'll want to know: \ru{журчать}}
\textbf{\ru{журчать}} (\emph{zhurchát}) — \textquotedblleft{}to babble, to murmur (of water)\textquotedblright{}.

\begin{quote}
\ru{Под} \ru{мостом} \ru{тихо} \ru{журчит} \ru{ручей}.
\end{quote}

\begin{quote}
\emph{Pod mostóm tíkho zhurchít ruchéy.}
\end{quote}

\begin{quote}
Under the bridge a stream babbles quietly.
\end{quote}

Imperfective. Present \emph{\ru{журчит}}, \emph{\ru{журчат}}.

\begin{cousinweb}
\emph{\ru{журчать}} is usually taken to imitate the sound of running water; EtymDB lists it without a source.
\end{cousinweb}

\begin{grammarlens}[title={What you've built}]
One more word for telling, explaining and describing at B1.
\end{grammarlens}

\subsection*{Your turn}
\begin{itemize}
\raggedright
  \item \emph{Say it:} \emph{zhurchát}
  \item \emph{Say it:} \emph{zhurchát}, once more
  \item \emph{Say it:} a sentence of your own with \emph{zhurchát}
  \item \emph{Recall:} say the Russian for a spring, then the Russian for a pond, then say \emph{zhurchát} again
\end{itemize}

\begin{tcolorbox}[breakable,colback=teal!4,colframe=teal!35!black,title={Before you move on}]
What does \ru{журчать} mean? (\textquotedblleft{}To babble, to murmur (of water)\textquotedblright{}.) Say it once more.
\end{tcolorbox}
